\chapter{超分子化学}\label{ch:b3:supramolecular}

1962 年，查尔斯·佩德森在寻找别的东西时分离出一种环状聚醚，
它能将高锰酸钾溶于苯中，并使溶液变成紫色。此环
已包裹住钾离子，把它的电荷隐藏于烃类表面之下，
并把高锰酸根作为抗衡离子一同带入苯中。
超越单个分子的化学，即由比共价键弱的作用力
结合而成的组装体的化学，正是从这一片紫色溶液开始的。
本章讨论这些作用力，
解释某些\omterm{def:b3:supramolecular:supramolecular}{主体}为何能如此精准地挑选\omterm{def:b3:supramolecular:supramolecular}{客体}，说明结合常数如何测定，
最后以分子被穿线、互锁成机器作结。

\begin{recall}
本科一年级分册讨论了分子间作用力、氢键、溶剂化、
疏水分子、配合物、生成常数和螯合物；本科二年级分册讨论了
化学势、吉布斯自由能和最小二乘拟合。
\cref{ch:b3:advanced-nmr}介绍了快速交换、慢速交换和合并现象，
\cref{ch:b3:rate-theories}介绍了\omterm{thm:b3:rate-theories:eyring}{艾林方程}，以及
\cref{ch:b3:partition-functions}介绍了气体的平动熵。
\end{recall}

\section{非共价相互作用}

\begin{definition}[超分子化学]\label{def:b3:supramolecular:supramolecular}
\emph{超分子化学}\index{超分子化学}是研究由非共价相互作用结合在一起的分子组装体的
化学。在\emph{主客体配合物}\index{主客体配合物}中，
\emph{主体}\index{主体}是较大的分子，带有空腔或裂缝，
结合位点汇聚其上，\emph{客体}\index{客体}是它所结合的
较小分子或离子。
\end{definition}

这些相互作用即本科一年级分册中的离子–偶极、偶极–偶极、
氢键和伦敦力，外加几种有专门名称的作用。

\begin{definition}[非共价相互作用]\label{def:b3:supramolecular:interactions}
\emph{$\pi$ 堆积}\index{$\pi$ 堆积}是平面芳香环之间的
面对面排列（通常错开或倾斜）的吸引作用。\emph{阳离子–$\pi$
相互作用}\index{阳离子–$\pi$ 相互作用}是阳离子与芳香环
富电子面之间的吸引作用。\emph{卤键}\index{卤键}是共价结合的
卤原子上与其键相对的缺电子端与
路易斯碱之间的吸引作用。\emph{疏水
效应}\index{疏水效应}是非极性表面在水中聚集到一起的
趋势，主要驱动力是非极性表面周围有序水分子的
释放。
\end{definition}

离子–偶极与阳离子–$\pi$ 相互作用最强，很大程度上取决于与离子竞争的
溶剂；氢键次之，
其后是 $\pi$ 堆积与\omterm{def:b3:supramolecular:interactions}{卤键}，再其后是伦敦力。伦敦力单个都很弱，
但在大接触面上累加起来便不容忽视。在水中，\omterm{def:b3:supramolecular:interactions}{疏水效应}往往占主导：
若某\omterm{def:b3:supramolecular:supramolecular}{主体}设计为在氯仿中靠氢键结合\omterm{def:b3:supramolecular:supramolecular}{客体}，它在水中可能
根本无法结合该\omterm{def:b3:supramolecular:supramolecular}{客体}，因为水自身的氢键会与之竞争。

\section{分子识别}

\begin{definition}[分子识别]\label{def:b3:supramolecular:recognition}
\emph{分子识别}\index{分子识别}是\omterm{def:b3:supramolecular:supramolecular}{主体}对\omterm{def:b3:supramolecular:supramolecular}{客体}的选择性
结合。它依赖于\emph{互补性}\index{互补性}，
即主\omterm{def:b3:supramolecular:supramolecular}{客体}在尺寸、形状和结合位点上的匹配，
并通过\emph{预组织}\index{预组织}得到加强，即\omterm{def:b3:supramolecular:supramolecular}{主体}的结合位点
在结合之前就已保持在配合物所需的几何构型中。
\end{definition}

结合常数 $K_a = [\mathrm{HG}]/([\mathrm H][\mathrm G])$ 在本科一年级分册的
意义上即为生成常数，并且 $\Delta G^\circ = -RT\ln K_a$。柔性\omterm{def:b3:supramolecular:supramolecular}{主体}在结合时必须冻结其
构象，这要付出熵的代价；刚性、已\omterm{def:b3:supramolecular:recognition}{预组织}的\omterm{def:b3:supramolecular:supramolecular}{主体}
则已在合成时付出了这一代价。

\begin{definition}[螯合效应与大环效应]\label{def:b3:supramolecular:macrocyclic}
\emph{螯合效应}\index{螯合效应}是指：含多齿配体的配合物
比含相同供体原子、数目相当的单齿配体的配合物更稳定。
\emph{大环
效应}\index{大环效应}是指：当
多齿配体为环状而非开链时，所得到的进一步稳定化。
\end{definition}

\begin{proposition}[螯合效应主要是熵效应]\label{prop:b3:supramolecular:chelate-entropy}
在 $\mathrm{[M(NH_3)_6]^{2+} + 3\,en \rightleftharpoons [M(en)_3]^{2+} + 6\,NH_3}$ 中，en 为乙烷-1,2-二胺，
反应两侧存在相同的六个金属–氮键，反应的驱动力是
自由分子数目的增加。
\end{proposition}

\begin{proof}[论证]
形成和断裂的键属于同一类，因此 $\Delta_rH^\circ$ 很小。独立分子的
数目由四个增至七个：整体分子多出三个平动自由度，
每个在溶液中每开尔文每摩尔可贡献数十焦耳的熵
（\cref{ch:b3:partition-functions}），因此 $\Delta_rS^\circ > 0$ 且 $K > 1$。
从动力学角度看，螯合配体的一端一旦结合，另一端就被以很高的有效浓度保持在
金属附近，若它松开便立即结合。
\end{proof}

\begin{proposition}[焓–熵补偿]\label{prop:b3:supramolecular:compensation}
在一系列相关的\omterm{def:b3:supramolecular:supramolecular}{主客体配合物}中，结合时较负的 $\Delta H^\circ$
通常伴随较负的 $\Delta S^\circ$，因此 $\Delta G^\circ$ 的变化小于其中任一项的变化（
这是一个实验观察结果）。
\end{proposition}

结合更紧的配合物结合得更强，但也更多地限制了\omterm{def:b3:supramolecular:supramolecular}{主体}和\omterm{def:b3:supramolecular:supramolecular}{客体}的
运动；与\omterm{def:b3:supramolecular:supramolecular}{客体}形成的更强氢键，往往意味着与被它取代的水之间的
氢键变弱。决定选择性的吉布斯自由能，是两个大项的
微小差值。

\begin{definition}[主体]\label{def:b3:supramolecular:hosts}
\emph{冠醚}\index{冠醚}是大环聚醚，例如
18-冠-6，$\ce{(CH2CH2O)6}$。\emph{穴醚}\index{穴醚}是双环
聚醚，带有两个氮桥头，能在三维空间中包围阳离子，
形成\emph{穴合物}\index{穴合物}。
\emph{环糊精}\index{环糊精}是由葡萄糖单元构成的环状低聚物，
呈截头圆锥形，内部为疏水空腔，
边缘带有羟基。\emph{杯芳烃}\index{杯芳烃}是杯状大环，
由酚单元经亚甲基桥连接而成。
\end{definition}

\begin{omfigure}
\begin{tikzpicture}[font=\small]
  % 18-crown-6 with K+
  \foreach \k in {0,...,17} {
    \pgfmathsetmacro\ang{90 + 20*\k}
    \pgfmathsetmacro\nxt{110 + 20*\k}
    \draw[black!70] (\ang:1.55) -- (\nxt:1.55);
  }
  \foreach \k in {0,3,...,15} {
    \pgfmathsetmacro\ang{90 + 20*\k}
    \node[fill=white, inner sep=0.5pt, text=omThm] (o\k) at (\ang:1.55) {O};
    \draw[omThm!60, densely dashed] (0,0) -- (\ang:1.32);
  }
  \node[circle, fill=omDef!25, draw=omDef, inner sep=1pt] at (0,0) {$\ce{K+}$};
  \node[font=\scriptsize] at (0,-2.1) {18-冠-6 与 $\ce{K+}$};
  % [2.2.2]cryptand with K+
  \begin{scope}[xshift=5.2cm]
    \node (nt) at (0,1.7) {N};
    \node (nb) at (0,-1.7) {N};
    \foreach \x/\n in {-1.25/a, 0/b, 1.25/c} {
      \node[text=omThm, fill=white, inner sep=0.5pt] (oa\n) at (\x,0.55) {O};
      \node[text=omThm, fill=white, inner sep=0.5pt] (ob\n) at (\x,-0.55) {O};
      \draw[black!70] (nt) to[out=-90+40*\x, in=90] (oa\n) -- (ob\n) to[out=-90, in=90-40*\x] (nb);
    }
    \node[circle, fill=omDef!25, draw=omDef, inner sep=1pt] at (0.62,0) {$\ce{K+}$};
    \node[font=\scriptsize] at (0,-2.4) {[2.2.2]穴醚：$\ce{K+}$ 穴合物};
  \end{scope}
\end{tikzpicture}

{\small 左：18-冠-6，由十二个碳原子（顶点）和六个氧原子构成的环，
环中心有一个钾离子，被六个氧原子配位。右：
[2.2.2]\omterm{def:b3:supramolecular:hosts}{穴醚}，两个氮桥头之间连有三条各含两个醚氧的链，
示意图：阳离子被包围在三维笼状空间之中，
由六个氧原子和两个氮原子包围。}
\end{omfigure}

\omterm{def:b3:supramolecular:hosts}{冠醚}按离子大小选择阳离子：18-冠-6 与 $\ce{K+}$ 的结合，比
较小的 $\ce{Na+}$ 强（后者无法同时接触六个氧原子），也比
较大的 $\ce{Cs+}$ 强（后者只能位于环的上方）。
\omterm{def:b3:supramolecular:hosts}{穴醚}\omterm{def:b3:supramolecular:recognition}{预组织}程度更高，能将离子完全包裹，结合
更强，选择性也更高。阳离子被包藏在\omterm{def:b3:supramolecular:supramolecular}{主体}内部，会把阴离子带入非极性溶剂中，
在其中阴离子未被溶剂化，活性很强，这正是相转移催化的基础。

\section{测定结合}

\begin{proposition}[快速交换]\label{prop:b3:supramolecular:fast-exchange}
当\omterm{def:b3:supramolecular:supramolecular}{主体}在游离态与结合态之间的交换速度远快于二者共振频率之差时，
其核磁共振信号出现在平均值
$\delta_{\mathrm{obs}} = x_{\mathrm{free}}\delta_{\mathrm{free}} + x_{\mathrm{bound}}\delta_{\mathrm{bound}}$ 处，其中 $x$ 为\omterm{def:b3:supramolecular:supramolecular}{主体}在各状态中的分数。
\end{proposition}

\begin{proof}
每个\omterm{def:b3:supramolecular:supramolecular}{主体}分子在测量过程中多次往返于两种状态之间；其进动频率
经时间平均后，是两者的加权平均，
权重为其在各状态中停留时间的分数。平衡时，这些分数即为各状态中分子的分数
（见\cref{ch:b3:advanced-nmr}）。
\end{proof}

\begin{theorem}[精确的 1:1 等温线]\label{thm:b3:supramolecular:isotherm}
对于 $\mathrm{H + G \rightleftharpoons HG}$，设总浓度分别为 $H_0$ 和 $G_0$，缔合常数为 $K$，则
\[
  [\mathrm{HG}] = \frac{1}{2}\Bigl[\Bigl(H_0 + G_0 + \frac1K\Bigr) - \sqrt{\Bigl(H_0 + G_0 + \frac1K\Bigr)^2 - 4H_0G_0}\,\Bigr].
\]
\end{theorem}

\begin{proof}
记 $c = [\mathrm{HG}]$，由 $K = c/((H_0 - c)(G_0 - c))$ 可得 $c^2 - (H_0 + G_0 + 1/K)c + H_0G_0 = 0$。两根中，
只有较小者同时小于 $H_0$ 和 $G_0$（两根之积为 $H_0G_0$，两根之和
超过 $H_0 + G_0$），即取负号的那一根。
\end{proof}

\begin{omfigure}
\begin{tikzpicture}
\begin{axis}[name=a, width=0.48\linewidth, height=5.5cm, xmin=0, xmax=6, ymin=0, ymax=1.05,
  axis lines=left, xlabel={客体的当量数，$G_0/H_0$}, ylabel={结合的主体分数},
  tick label style={font=\scriptsize}, label style={font=\small},
  ]
\addplot[omMeth, thick] table[x=equiv, y=K4] {figdata/out/bachelor-3-binding-titration.dat};
\addplot[omDef, thick] table[x=equiv, y=K3] {figdata/out/bachelor-3-binding-titration.dat};
\addplot[omThm, thick] table[x=equiv, y=K2] {figdata/out/bachelor-3-binding-titration.dat};
\node[font=\scriptsize, omMeth, anchor=north] at (axis cs:4,0.94) {$K = 10^4$};
\node[font=\scriptsize, omDef, anchor=north] at (axis cs:4,0.68) {$K = 10^3$};
\node[font=\scriptsize, omThm, anchor=north] at (axis cs:4,0.27) {$K = 10^2$};
\end{axis}
\begin{axis}[at={(a.east)}, anchor=west, xshift=1.4cm, width=0.48\linewidth, height=5.5cm,
  xmin=0, xmax=1, ymin=0, ymax=1.05, axis lines=left,
  xlabel={主体的摩尔分数}, ylabel={$[\text{复合物}]/c_{\text{总}}$},
  tick label style={font=\scriptsize}, label style={font=\small}]
\addplot[omDef, thick] table[x=x, y=HG] {figdata/out/bachelor-3-binding-titration-b.dat};
\addplot[omThm, thick] table[x=x, y=HG2] {figdata/out/bachelor-3-binding-titration-b.dat};
\draw[black!35, dashed] (axis cs:0.5,0) -- (axis cs:0.5,0.55) (axis cs:0.3333,0) -- (axis cs:0.3333,0.38);
\node[font=\scriptsize, omDef, anchor=south] at (axis cs:0.5,0.52) {HG};
\node[font=\scriptsize, omThm, anchor=south] at (axis cs:0.25,0.36) {$\mathrm{HG_2}$};
\end{axis}
\end{tikzpicture}

{\small 左：在 $H_0 = \qty{1.0}{mmol/L}$ 的滴定中，三种缔合常数对应的
结合\omterm{def:b3:supramolecular:supramolecular}{主体}分数（缔合常数单位为 $\unit{L.mol^{-1}}$）：$K$ 越大，在一当量处的转折越陡；
当 $KH_0 \gg 1$ 时，曲线只能说明 $K$ 很大。右：\omterm{def:b3:supramolecular:job}{乔布图}，即
在\omterm{def:b3:supramolecular:supramolecular}{主体}与\omterm{def:b3:supramolecular:supramolecular}{客体}总浓度保持恒定时配合物的浓度，
对 1:1 配合物在\omterm{def:b3:supramolecular:supramolecular}{主体}分数为 1/2 处取峰值，对 1:2 配合物则在 1/3 处取峰值（模型
常数）。}
\end{omfigure}

\begin{method}[由核磁共振滴定求结合常数]\label{met:b3:supramolecular:nmr-titration}
\begin{enumerate}
  \item 保持\omterm{def:b3:supramolecular:supramolecular}{主体}浓度 $H_0$ 不变，逐次加入\omterm{def:b3:supramolecular:supramolecular}{客体}，\omterm{def:b3:supramolecular:supramolecular}{客体}的当量数从 0 增至数个
        当量；每加一次都记录\omterm{def:b3:supramolecular:supramolecular}{主体}的一个信号（快速交换）。
  \item 模型：$\delta = \delta_{\mathrm{free}} + (\delta_{\mathrm{bound}} - \delta_{\mathrm{free}})[\mathrm{HG}]/H_0$，其中 $[\mathrm{HG}]$ 由精确等温线给出。
  \item 用非线性最小二乘法拟合 $K$ 与 $\delta_{\mathrm{bound}}$，使残差平方和
        最小；标准误差取自该和的曲率，
        （即协方差矩阵）。
  \item 检查残差是否存在趋势（化学计量比可能有误），并检查
        结合分数是否大致覆盖 20\,\% 至 80\,\%（此时数据方能约束 $K$）。
\end{enumerate}
\end{method}

\begin{definition}[乔布图]\label{def:b3:supramolecular:job}
\emph{乔布图}\index{乔布图}（连续变化法）是指
在两组分总浓度保持不变时，配合物的浓度（或与之成正比的信号）
对其中一组分的摩尔分数作图所得的
曲线。
\end{definition}

\begin{proposition}[乔布图的最大值]\label{prop:b3:supramolecular:job}
对于单一配合物 $\mathrm{H_nG_m}$，\omterm{def:b3:supramolecular:job}{乔布图}在\omterm{def:b3:supramolecular:supramolecular}{主体}的摩尔分数
$x = n/(n + m)$ 处取得最大值，与结合常数无关。
\end{proposition}

\begin{proof}
设总浓度为 $T$，$H_t = xT$，$G_t = (1 - x)T$，配合物浓度为 $c$，游离态 $[\mathrm H] = H_t - nc$，
$[\mathrm G] = G_t - mc$，且 $c = \beta[\mathrm H]^n[\mathrm G]^m$。在取得最大值处对 $x$ 求导，此时
$\dd c/\dd x = 0$：$0 = \beta\bigl(n[\mathrm H]^{n-1}[\mathrm G]^m T - m[\mathrm H]^n[\mathrm G]^{m-1}T\bigr)$，故 $n[\mathrm G] = m[\mathrm H]$。代入得
$n(G_t - mc) = m(H_t - nc)$，于是 $nG_t = mH_t$，即 $n(1 - x) = mx$：$x = n/(n + m)$。
\end{proof}

\begin{method}[乔布图的测定]\label{met:b3:supramolecular:job}
\begin{enumerate}
  \item 配制\omterm{def:b3:supramolecular:supramolecular}{主体}与\omterm{def:b3:supramolecular:supramolecular}{客体}的等物质的量浓度母液。
  \item 在总体积保持不变的条件下，按 $x$ 从 0 到 1 的比例混合两种母液。
  \item 测量与配合物成正比的性质（经游离物种校正的吸光度，
        或核磁共振中的 $x\,\Delta\delta$）。
  \item 由最大值所在位置读出化学计量比。
\end{enumerate}
\end{method}

\begin{definition}[等温滴定量热法]\label{def:b3:supramolecular:itc}
\emph{等温滴定量热法}\index{等温滴定量热法}
（ITC）用于测量在恒温下，每次向\omterm{def:b3:supramolecular:supramolecular}{主体}溶液中注入\omterm{def:b3:supramolecular:supramolecular}{客体}溶液时
所释放或吸收的热量。
\end{definition}

\begin{proposition}[每次注入的热量]\label{prop:b3:supramolecular:itc}
忽略稀释效应时，第 $i$ 次注入样品池（体积为 $V_0$）所放出的热量为
$q_i = \Delta H^\circ V_0\bigl([\mathrm{HG}]_i - [\mathrm{HG}]_{i-1}\bigr)$。
\end{proposition}

\begin{proof}
热量等于反应焓乘以注入过程中生成的配合物的物质的量，
该物质的量等于 $[\mathrm{HG}]$ 的变化量乘以样品池体积。
\end{proof}

\begin{method}[读取 ITC 等温线]\label{met:b3:supramolecular:itc}
\begin{enumerate}
  \item 以每摩尔注入\omterm{def:b3:supramolecular:supramolecular}{客体}的热量对摩尔比作图。
  \item 饱和前平台的高度给出 $\Delta H^\circ$；拐点处的摩尔比
        给出化学计量比；曲线的陡峭程度给出 $K$。
  \item 用精确等温线同时拟合这三个量；然后由 $\Delta G^\circ = -RT\ln K$ 和
        $T\Delta S^\circ = \Delta H^\circ - \Delta G^\circ$ 求出各热力学量。
\end{enumerate}
\end{method}

因此，单次 ITC 实验即可给出结合过程完整的热力学特征，
即焓与熵，而核磁共振或紫外–可见滴定则只能
通过一系列温度下的测量才能得到。

\begin{inthelab}[核磁共振滴定]
将\omterm{def:b3:supramolecular:supramolecular}{主体}溶于氘代溶剂中，配成约 1 毫摩尔每升的溶液，
置于核磁管中并记录其谱图。取\omterm{def:b3:supramolecular:supramolecular}{客体}的浓溶液，
该溶液以\omterm{def:b3:supramolecular:supramolecular}{主体}溶液配制，使\omterm{def:b3:supramolecular:supramolecular}{主体}浓度保持不变，
分小份用微量注射器加入；每次加入后
将核磁管摇匀，在探头温度下平衡，并
记录谱图。逐点读取发生移动的\omterm{def:b3:supramolecular:supramolecular}{主体}信号的化学位移，
并进行拟合。
\end{inthelab}

\section{主体}

\omterm{def:b3:supramolecular:hosts}{环糊精}由淀粉经酶作用制得，由六个、七个或八个葡萄糖单元构成
（分别记为 $\alpha$、$\beta$、$\gamma$）。其疏水空腔可在
水中容纳\omterm{def:b3:supramolecular:supramolecular}{客体}的非极性部分：它们能增溶药物、负载香料，并从
细胞膜中除去胆固醇。\omterm{def:b3:supramolecular:hosts}{杯芳烃}与形如南瓜的葫芦脲二者均为刚性
大环，由甘脲单元构成、两端口缘带有羰基，能与阳离子
及有机阳离子在水中发生非常强的结合。

\begin{omfigure}
\begin{tikzpicture}[font=\small]
  \draw[black!70, fill=omDef!8] (-1.9,1.2) -- (-1.2,-1.2) arc[start angle=180, end angle=360, x radius=1.2, y radius=0.3]
    -- (1.9,1.2);
  \draw[black!70, fill=omDef!18] (0,1.2) ellipse[x radius=1.9, y radius=0.45];
  \draw[black!40, dashed] (-1.2,-1.2) arc[start angle=180, end angle=0, x radius=1.2, y radius=0.3];
  \node[draw=omThm, thick, regular polygon, regular polygon sides=6, minimum size=8mm, inner sep=0pt] at (0,0.2) {};
  \draw[omThm, thick] (0,0.62) -- (0,2.0);
  \node[font=\scriptsize, omThm, anchor=west] at (0.1,2.0) {客体};
  \node[font=\scriptsize, anchor=west] at (2.0,1.2) {宽口缘：仲羟基};
  \node[font=\scriptsize, anchor=west] at (1.3,-1.25) {窄口缘：伯羟基};
  \node[font=\scriptsize, anchor=east] at (-1.75,0) {疏水空腔};
\end{tikzpicture}

{\small \omterm{def:b3:supramolecular:hosts}{环糊精}呈其所形成的截头锥体，其空腔中包结一个芳香
\omterm{def:b3:supramolecular:supramolecular}{客体}。两端口缘上的羟基使外表面
可溶于水；内壁衬有 C--H 基团与糖苷氧原子，
为非极性。}
\end{omfigure}

\section{自组装与分子机器}

\begin{definition}[自组装]\label{def:b3:supramolecular:self-assembly}
\emph{自组装}\index{自组装}是指组分在其自身所含信息的控制下，
自发且可逆地组织成确定结构的过程，此信息即
组分自身的形状与
结合位点。
\end{definition}

具有固定配位角的金属离子与刚性桥联配体，
可一步组装成正方形、笼状与胶囊状结构，因为每个键都可断开
与重新形成，直至达到最稳定的结构，
且无张力、无悬挂位点。DNA 中的碱基配对与蛋白质折叠，
则是同一思想在更大尺度上的体现。

\begin{definition}[机械互锁分子]\label{def:b3:supramolecular:interlocked}
称\emph{机械互锁分子}\index{机械互锁分子}为一类分子，其各部分彼此之间
无化学键相连，但若不断开共价键，它们便无法分开。
\emph{轮烷}\index{轮烷}是一个环穿套在轴上的结构，
轴的两端带有庞大的端基；
\emph{索烃}\index{索烃}由相互套合的环组成。
\end{definition}

\begin{definition}[模板合成]\label{def:b3:supramolecular:template}
\emph{模板合成}\index{模板合成}利用一种暂时性的相互作用，
常与金属离子有关，使活性组分保持
反应所需的几何构型，从而完成一个本来不易发生的反应，
例如一个环围绕另一个环闭合。
\end{definition}

\begin{omfigure}
\begin{tikzpicture}[font=\small, >=Stealth]
  % step 1: two crescents around Cu+
  \draw[omDef, line width=2pt] (0.3,0) arc[start angle=0, end angle=300, radius=0.75];
  \draw[omThm, line width=2pt] (-0.3,0) arc[start angle=180, end angle=-120, radius=0.75];
  \fill[omProp] (0,0) circle (0.13);
  \node[font=\scriptsize] at (0,-1.35) {$\ce{Cu+}$ 与两个};
  \node[font=\scriptsize] at (0,-1.7) {交叉的开链配体};
  \draw[->, thick] (1.4,0) -- (2.3,0) node[midway, above, font=\scriptsize] {闭合};
  % step 2: closed interlocked rings with Cu
  \begin{scope}[xshift=3.6cm]
    \draw[omDef, line width=2pt] (-0.35,0) circle (0.75);
    \draw[omThm, line width=2pt] (0.35,0) circle (0.75);
    \draw[white, line width=5pt] ([shift={(52:0.75)}]-0.35,0) arc[start angle=52, end angle=72, radius=0.75];
    \draw[omDef, line width=2pt] ([shift={(48:0.75)}]-0.35,0) arc[start angle=48, end angle=76, radius=0.75];
    \fill[omProp] (0,0) circle (0.13);
    \node[font=\scriptsize] at (0,-1.35) {套合};
  \end{scope}
  \draw[->, thick] (5.0,0) -- (5.9,0) node[midway, above, font=\scriptsize] {$-\ce{Cu+}$};
  \begin{scope}[xshift=7.2cm]
    \draw[omDef, line width=2pt] (-0.35,0) circle (0.75);
    \draw[omThm, line width=2pt] (0.35,0) circle (0.75);
    \draw[white, line width=5pt] ([shift={(52:0.75)}]-0.35,0) arc[start angle=52, end angle=72, radius=0.75];
    \draw[omDef, line width=2pt] ([shift={(48:0.75)}]-0.35,0) arc[start angle=48, end angle=76, radius=0.75];
    \node[font=\scriptsize] at (0,-1.35) {[2]索烃};
  \end{scope}
  % rotaxane
  \begin{scope}[xshift=10.2cm]
    \draw[black!70, line width=1.6pt] (-1.1,0) -- (1.1,0);
    \fill[black!60] (-1.25,0) circle (0.2) (1.25,0) circle (0.2);
    \draw[omThm, line width=2pt] (0,0) ellipse[x radius=0.18, y radius=0.55];
    \node[font=\scriptsize] at (0,-1.35) {[2]轮烷};
  \end{scope}
\end{tikzpicture}

{\small \omterm{def:b3:supramolecular:interlocked}{索烃}的\omterm{def:b3:supramolecular:template}{模板合成}（示意）：一个铜（I）离子托住两个
菲咯啉类配体，二者交叉成直角，位于其四面体配位之中；
将每个配体闭合成环，在金属周围得到两个相互套合的环；
移去铜即得到游离的\omterm{def:b3:supramolecular:interlocked}{索烃}。右图为\omterm{def:b3:supramolecular:interlocked}{轮烷}，
即一个环被两个端基套在轴上。}
\end{omfigure}

\begin{definition}[分子机器]\label{def:b3:supramolecular:machine}
\emph{分子机器}\index{分子机器}是一种组装体，其组分
能以可控方式相对运动，以响应一种外界刺激
（光、氧化还原变化、pH 变化），并能对外做功。
\end{definition}

在分子穿梭器中，\omterm{def:b3:supramolecular:interlocked}{轮烷}轴上的一个环停留在两个位点之一；
刺激使其与第一个位点的结合减弱，环便移至第二个位点，
反向刺激则使其返回。基于拥挤烯烃的光驱动旋转马达，
只沿一个方向转动，其交替进行的是光化学异构化
与不能逆向进行的热步骤。穿梭器两个位点之间的穿梭速率，
可由变温核磁共振测定，
方法是由两种形式信号的合并（见\cref{ch:b3:advanced-nmr}）。

\begin{safety}{\ghs[0.62cm]{GHS07}}
% ledger: ghs4:18C6
18-冠-6：吞食有害，刺激皮肤、眼睛和呼吸道。
\omterm{def:b3:supramolecular:hosts}{冠醚}会将金属盐带入有机溶剂，并可穿透皮肤；
称量和操作时须戴手套。
\end{safety}

\begin{history}[超分子化学及其机器]
查尔斯·佩德森的\omterm{def:b3:supramolecular:hosts}{冠醚}（1967 年）、让-马里·莱恩的\omterm{def:b3:supramolecular:hosts}{穴醚}
以及唐纳德·克拉姆的\omterm{def:b3:supramolecular:recognition}{预组织}\omterm{def:b3:supramolecular:supramolecular}{主体}奠定了该领域的基础；三人于 1987 年
共同获得诺贝尔化学奖。让-皮埃尔·萨瓦热于 1983 年以铜模板法
高产率地制得\omterm{def:b3:supramolecular:interlocked}{索烃}，弗雷泽·斯托达特构建了\omterm{def:b3:supramolecular:interlocked}{轮烷}穿梭器，
本·费林加于 1999 年制成光驱动旋转马达；三人因在\omterm{def:b3:supramolecular:machine}{分子机器}
的设计与合成方面的成就，共同获得 2016 年诺贝尔化学奖。
\end{history}

\section{练习}

\begin{exercise}[$\star$]\label{exo:b3:supramolecular:1}
指出下列各体系中的主要相互作用：$\ce{K+}$ 在 18-冠-6 中；两个苯环面对面；
$\ce{K+}$ 位于苯环上方；五氟碘苯与吡啶；鸟嘌呤–胞嘧啶
碱基对。
\end{exercise}

\begin{exercise}[$\star$]\label{exo:b3:supramolecular:2}
计算\omterm{def:b3:supramolecular:supramolecular}{主客体配合物}在 \qty{298}{K} 时的 $\Delta G^\circ$，已知 $K_a = \qty{1.0e4}{L.mol^{-1}}$。
\end{exercise}

\begin{exercise}[$\star$]\label{exo:b3:supramolecular:3}
在甲醇中，$\log K = 6.1$ 对应 $\ce{K+}$，4.4 对应 $\ce{Na+}$，均与 18-冠-6 结合（数据见本题）。
计算选择性，以及 $\Delta G^\circ$ 的差值。
\end{exercise}

\begin{exercise}[$\star$]\label{exo:b3:supramolecular:4}
\omterm{def:b3:supramolecular:job}{乔布图}在\omterm{def:b3:supramolecular:supramolecular}{主体}摩尔分数为 0.33 处出现峰值，求其化学计量比。
\end{exercise}

\begin{exercise}[$\star\star$]\label{exo:b3:supramolecular:5}
\omterm{def:b3:supramolecular:supramolecular}{主体}浓度为 \qty{2.0}{mmol/L}，$K = \qty{1500}{L.mol^{-1}}$，$\delta_{\mathrm{free}} = \qty{3.560}{ppm}$ 和 $\delta_{\mathrm{bound}} = \qty{3.720}{ppm}$，
与一当量的\omterm{def:b3:supramolecular:supramolecular}{客体}结合。计算结合分数以及观测位移。
\end{exercise}

\begin{exercise}[$\star\star$]\label{exo:b3:supramolecular:6}
对于镍(II)，与氨的 $\log\beta_6 = 8.6$，与乙烷-1,2-二胺的 $\log\beta_3 = 18.3$（本题
数据）。计算下列反应的平衡常数，以及 \qty{298}{K} 时的 $\Delta G^\circ$：
$\ce{[Ni(NH3)6]^2+} + 3\,\mathrm{en} \rightleftharpoons \ce{[Ni(en)3]^2+} + \ce{6NH3}$，并说明驱动该反应的因素。
\end{exercise}

\begin{exercise}[$\star\star$]\label{exo:b3:supramolecular:7}
一次等温滴定量热实验在 \qty{298}{K} 时得到 $K = \qty{2.0e5}{L.mol^{-1}}$，$\Delta H^\circ = \qty{-40}{kJ/mol}$。求 $\Delta G^\circ$、$T\Delta S^\circ$ 与
$\Delta S^\circ$。
\end{exercise}

\begin{exercise}[$\star\star$]\label{exo:b3:supramolecular:8}
说明四面体构型的铜(I)及其构型在索瓦日的
\omterm{def:b3:supramolecular:interlocked}{索烃}合成中所起的作用。反应结束后，如何除去铜？
\end{exercise}

\begin{exercise}[$\star\star$]\label{exo:b3:supramolecular:9}
\omterm{def:b3:supramolecular:interlocked}{轮烷}穿梭体的两个信号相距 \qty{120}{Hz}，在 \qty{300}{K} 时
合并（本题数据）。计算合并时的交换速率常数，以及
穿梭过程的\omterm{def:b3:rate-theories:activation}{活化吉布斯自由能}。
\end{exercise}

\begin{exercise}[$\star\star\star$]\label{exo:b3:supramolecular:10}
某药物的本征溶解度为 \qty{0.10}{mmol/L}，与一种\omterm{def:b3:supramolecular:hosts}{环糊精}形成 1:1 配合物，
结合常数 $K = \qty{500}{L.mol^{-1}}$（本题数据）。证明：在总\omterm{def:b3:supramolecular:hosts}{环糊精}浓度为 $C_t$ 的
溶液中，总溶解度为 $S_0 + KS_0C_t/(1 + KS_0)$，并计算 $C_t = \qty{10}{mmol/L}$ 时的值。
\end{exercise}

\begin{exercise}[$\star\star\star$]\label{exo:b3:supramolecular:11}
由精确等温线证明，当 $KH_0 \gg 1$ 时，结合分数先等于
当量数（直至 1），此后恒为 1。为什么这样的滴定只能给出
$K$ 的下限？
\end{exercise}

\begin{exercise}[$\star\star\star$]\label{exo:b3:supramolecular:12}
对于\cref{exo:b3:supramolecular:5}中的\omterm{def:b3:supramolecular:supramolecular}{主体}，假设缔合过程为扩散控制，
即 $k_{\mathrm{on}} = \qty{1e9}{L.mol^{-1}.s^{-1}}$。计算 $k_{\mathrm{off}}$，并证明
在 \qty{400}{MHz} 的核磁共振时间尺度上交换较快。
\end{exercise}

\section{周末问题：紫色苯}

\begin{problem}[{周末问题——紫色苯：冠醚中的钾，
高锰酸根向有机溶剂中的萃取，
由 NMR 滴定得到的结合常数，以及为何只需催化量的冠醚即可}]
\label{pb:b3:supramolecular:1}
% ledger: const:R, ghs:KMnO4, ghs:toluene, rion:K+VI, rion:Na+VI
题目数据。在 \qty{298}{K} 的甲醇中，与 18-冠-6 配位时，$\log K = 6.1$（$\ce{K+}$），4.4（$\ce{Na+}$）。
萃取：水相中含 \qty{0.10}{mol/L} $\ce{KCl}$ 与 \qty{1.0}{mmol/L} $\ce{KMnO4}$，
用等体积含\omterm{def:b3:supramolecular:hosts}{冠醚} L 的甲苯振荡萃取；唯一被萃取的
物种为离子对 $\ce{[KL]+}\ce{MnO4-}$，其中
$K_{\mathrm{ex}} = [\ce{[KL]+MnO4-}]_{\mathrm{org}}/([\ce{K+}]_{\mathrm{aq}}[\ce{MnO4-}]_{\mathrm{aq}}[\mathrm L]_{\mathrm{org}}) = \qty{1.0e5}{L^2.mol^{-2}}$，且\omterm{def:b3:supramolecular:hosts}{冠醚}留在甲苯中。对\omterm{def:b3:supramolecular:hosts}{冠醚}进行 NMR 滴定，
\omterm{def:b3:supramolecular:hosts}{冠醚}浓度为 \qty{2.0}{mmol/L}，滴定剂为氘代甲醇中的钠盐，记录 $\ce{CH2}$ 信号
的化学位移（ppm）随 $\ce{Na+}$ 当量数的变化：

\begin{center}\small
\begin{tabular}{l*{11}{c}}
\hline
当量数 & 0 & 0.25 & 0.5 & 0.75 & 1.0 & 1.5 & 2.0 & 3.0 & 4.0 & 6.0 & 8.0\\
$\delta$ & 3.560 & 3.590 & 3.612 & 3.635 & 3.649 & 3.674 & 3.688 & 3.697 & 3.704 & 3.708 & 3.714\\
\hline
\end{tabular}
\end{center}
% table: binding-titration-c

\medskip\noindent\textbf{第一部分——配合物。}
\begin{enumerate}
  \item 画出 18-冠-6 及其与钾形成的配合物。
  \item 计算 $\Delta G^\circ$，分别对应 $\ce{K+}$ 与 $\ce{Na+}$ 的结合。
  \item 计算 $\ce{K+}/\ce{Na+}$ 的选择性。
  \item 利用六配位半径 $r(\ce{K+}) = \qty{138}{pm}$ 与 $r(\ce{Na+}) = \qty{102}{pm}$，解释该选择性。
  \item 为什么该配合物可溶于甲苯，而 $\ce{KCl}$ 不溶？
  \item 为什么\omterm{def:b3:supramolecular:hosts}{冠醚}在水中的结合比在甲醇中弱？
\end{enumerate}

\medskip\noindent\textbf{第二部分——萃取。}
\begin{enumerate}[resume]
  \item 写出萃取平衡式。
  \item 设 $y$ 为萃取浓度，写出 $y$ 所满足的方程，当
        $[\mathrm L]_0 = \qty{1.0}{mmol/L}$ 时的方程。
  \item 解出该方程，并给出高锰酸根的萃取分数。
  \item 计算 $[\mathrm L]_0 = \qty{10}{mmol/L}$ 时的萃取分数。
  \item 计算 $[\mathrm L]_0 = \qty{0.10}{mmol/L}$ 时的萃取分数。
  \item 为什么甲苯呈紫色，以及为什么被萃取的
        高锰酸根是比在水中更强的氧化剂？
\end{enumerate}

\medskip\noindent\textbf{第三部分——一次 NMR 滴定。}
\begin{enumerate}[resume]
  \item 为什么移动的是一个平均化信号，而不是两个信号？
  \item 写出 $\delta$ 随\omterm{def:b3:supramolecular:supramolecular}{客体}浓度变化的模型。
  \item 用最小二乘法拟合 $K$ 和 $\delta_{\mathrm{bound}}$，给出 $K$ 及其标准误差。
  \item 给出残差的离散程度并加以评论。
  \item 结合分数在 20 至 80\,\% 之间时，当量数的范围为多少？
  \item 为什么同样的滴定无法测得钾的 $K$（$\log K = 6.1$）？
\end{enumerate}

\medskip\noindent\textbf{第四部分——相转移催化。}
\begin{enumerate}[resume]
  \item 甲苯中的烯烃被萃取出的高锰酸根氧化，该反应
        发生在何处？
  \item 反应后\omterm{def:b3:supramolecular:hosts}{冠醚}发生了什么变化，以及为什么
        只需催化量即可？
  \item 水相中过量的 $\ce{KCl}$ 有什么作用？
  \item 工业上哪些较廉价的催化剂可以起同样的作用？
  \item 为什么如今用甲苯，而佩德森当年用的是苯？
  \item 写出结果，即在
        $[\mathrm L]_0 = \qty{1.0}{mmol/L}$ 时高锰酸根的萃取分数。
\end{enumerate}
\end{problem}
